
\documentclass[lettersize,journal]{IEEEtran}
\IEEEoverridecommandlockouts

% -------------------- Essential Packages --------------------
\usepackage{amsmath,amsfonts}
\usepackage{amssymb,amsthm,mathtools,bm}
\usepackage{algorithm}
\usepackage{algorithmic}
\usepackage{graphicx}
\usepackage{textcomp}
\usepackage{xcolor}
\usepackage{cite}
\usepackage{url}
\usepackage{booktabs}
\usepackage{multirow}
\usepackage{array}
\usepackage{tikz}
\usepackage{pgfplots}
\pgfplotsset{compat=1.17}
\usetikzlibrary{positioning,arrows,shapes,calc,patterns,decorations.pathreplacing,backgrounds,fit}
\usepackage{soul}
\usepackage{enumitem}
\usepackage{listings}
\usepackage{microtype}
\usepackage[caption=false,font=footnotesize]{subfig}
\usepackage[hidelinks]{hyperref}
\usepackage[nameinlink,capitalise,noabbrev]{cleveref}
\usepackage{tabularx}
\usepgfplotslibrary{colormaps}
\usetikzlibrary{pgfplots.groupplots}

% -------------------- Custom Commands --------------------
\newcommand{\R}{\mathbb{R}}
\newcommand{\E}{\mathbb{E}}
\newcommand{\N}{\mathcal{N}}
\newcommand{\F}{\mathcal{F}}
\newcommand{\G}{\mathcal{G}}
\newcommand{\D}{\mathcal{D}}
\newcommand{\A}{\mathcal{A}}
\newcommand{\X}{\mathcal{X}}
\newcommand{\Y}{\mathcal{Y}}
\newcommand{\Z}{\mathcal{Z}}
\newcommand{\s}{\mathcal{S}}
\newcommand{\Tcal}{\mathcal{T}}
\newcommand{\Kcal}{\mathcal{K}}
\newcommand{\Mcal}{\mathcal{M}}
\newcommand{\Q}{\mathcal{Q}}
\newcommand{\U}{\mathcal{U}}
\newcommand{\V}{\mathcal{V}}
\newcommand{\W}{\mathcal{W}}
\newcommand{\h}{\mathcal{H}}
\newcommand{\softplus}{\operatorname{softplus}}
\newcommand{\Softmax}{\operatorname{Softmax}}
\DeclareMathOperator*{\argmin}{arg\,min}
\DeclareMathOperator*{\argmax}{arg\,max}
\DeclareMathOperator{\Tr}{Tr}
\DeclareMathOperator{\diag}{diag}
\DeclareMathOperator{\rank}{rank}
\DeclareMathOperator{\vecop}{vec}
\DeclareMathOperator{\prox}{prox}
\DeclareMathOperator{\KL}{KL}
\DeclareMathOperator{\ELBO}{ELBO}

% -------------------- Theorem Environments --------------------
\newtheorem{theorem}{Theorem}
\newtheorem{lemma}[theorem]{Lemma}
\newtheorem{proposition}[theorem]{Proposition}
\newtheorem{corollary}[theorem]{Corollary}
\newtheorem{definition}{Definition}
\newtheorem{assumption}{Assumption}
\newtheorem{remark}{Remark}

% Tighten spacing around floats
\setlength{\textfloatsep}{5pt}
\setlength{\floatsep}{5pt}
\setlength{\intextsep}{5pt}
\setlength{\skip\footins}{8pt}

% -------------------- Title & Authors --------------------
\title{CT-TGNN: Continuous-Time Temporal Graph Neural Networks for Real-Time Threat Propagation Analysis in Encrypted Zero-Trust Microservices}

\author{Roger~Nick~Anaedevha,~\IEEEmembership{Member,~IEEE},
        Alexander~Gennadevich~Trofimov,
        and~Yuri~Vladimirovich~Borodachev%
\thanks{R. N. Anaedevha and A. G. Trofimov are with the National Research Nuclear University MEPhI (Moscow Engineering Physics Institute), Moscow 115409, Russia.}%
\thanks{Y. V. Borodachev is with the Artificial Intelligence Research Center, National Research Nuclear University MEPhI, Moscow 115409, Russia.}%
\thanks{Corresponding email: ar006@campus.mephi.ru}%
\thanks{Manuscript received December 30, 2025.}}

% The paper headers
\markboth{IEEE Transactions on Services Computing,~Vol.~XX, No.~X, Month~2025}%
{Anaedevha \MakeLowercase{\textit{et al.}}: CT-TGNN for Real-Time Threat Propagation Analysis}

% -------------------- Document --------------------
\begin{document}

% Create title
\maketitle

\begin{abstract}
Zero-trust architectures require continuous verification of all network communications, yet the proliferation of encrypted traffic - now exceeding 87\% of enterprise communications—renders traditional deep packet inspection ineffective. Existing temporal graph neural networks operate on discrete time snapshots, missing critical inter-event attack propagation dynamics, while current encrypted traffic analyzers ignore network topology essential for detecting lateral movement and coordinated threats. This paper introduces Continuous-Time Temporal Graph Neural Networks (CT-TGNN) for encrypted traffic analysis in zero-trust microservices, enabling real-time threat propagation tracking through graph-structured continuous dynamics. We develop a unified framework coupling Neural Ordinary Differential Equations with temporal graph convolutions, allowing continuous evolution of node embeddings that capture both smooth state transitions and discrete security events across encrypted communication graphs. The proposed architecture integrates encrypted edge feature encoding that extracts timing patterns and protocol metadata without payload inspection, multi-scale temporal modeling spanning microseconds to hours, and zero trust policy enforcement through continuous authentication verification. Theoretical analysis establishes gradient stability for continuous graph dynamics via Lyapunov methods and proves convergence guarantees under dynamic graph topology shifts. Experimental validation on six diverse datasets—including microservices traces with 15 million encrypted API calls, the IoT-23 dataset with 325~GB of encrypted traffic, UNSW-NB15 temporal splits, CIC-IDS2018, NF-ToN-IoT, and NF-BoT-IoT demonstrates up to 98.3\% detection accuracy with 47~ms median latency, outperforming discrete temporal graph networks by up to 12.7\% in zero day lateral movement detection. The framework processes 8.7 million encrypted events per second while preserving privacy, enabling scalable and privacy-preserving threat intelligence across zero-trust service boundaries.
\end{abstract}

\begin{IEEEkeywords}
Services computing, zero-trust security, temporal graph neural networks, neural ordinary differential equations, encrypted traffic analysis, microservices security, continuous-time dynamics, real-time intrusion detection.
\end{IEEEkeywords}

% =========================================================
\section{Introduction}
\label{sec:introduction}

\IEEEPARstart{T}{hree} converging trends challenge network intrusion detection: zero-trust architectures that require continuous verification across microservices, encryption that protects 87\% of enterprise traffic, rendering deep packet inspection ineffective, and advanced persistent threats that exploit temporal patterns in encrypted communications through gradual reconnaissance over hours. Traditional systems sampling at fixed intervals miss attacks designed to evade periodic observation, while discrete temporal graph neural networks operating on hourly snapshots fail to capture continuous threat evolution between observations. Encrypted traffic analyzers that treat flows independently cannot detect coordinated exploration patterns across the topology of the service mesh.

This research introduces continuous-time temporal graph neural networks (CT-TGNN) for encrypted traffic analysis in zero-trust microservices, unifying Neural Ordinary Differential Equations~\cite{chen2018neural} for modeling security state evolution at arbitrary temporal resolution, temporal graph neural networks for capturing attack propagation through network topology, and real-time zero-trust policy enforcement through predicted security state verification. Building on recent advances including Temporal Adaptive Batch Normalization~\cite{salvi2024tabn} and transformer-based temporal point processes~\cite{zuo2020transformer, aandt2025transformer}, we formulate continuous-time graph dynamics through coupled ODEs with Lipschitz continuity conditions ensuring stable gradient flow, develop encrypted edge feature encoding extracting timing patterns from TLS 1.3 and QUIC traffic achieving 94.7\% attack classification, introduce multi-scale temporal convolutions spanning microseconds to hours, and provide theoretical convergence guarantees under dynamic graph topology shifts.

Comprehensive validation on six datasets demonstrates practical effectiveness. Evaluation of production Kubernetes traces (847 services, 15 million encrypted API calls) achieves 98.3\% detection accuracy with a median latency of 47 ms suitable for real-time enforcement. Experiments on IoT-23~\cite{garcia2020iot23} (325 GB encrypted traffic) demonstrate 96.8\% detection using only TLS timing patterns, while comparison with discrete approaches on UNSW-NB15~\cite{moustafa2015unsw} shows 12.7\% improvement on zero-day lateral movement. Validation on CIC-IDS2018~\cite{sharafaldin2018cic}, NF-ToN-IoT~\cite{sarhan2022nf}, and NF-BoT-IoT~\cite{sarhan2021netflow} confirms generalization across attack vectors, with scalability analysis confirming 8.7 million encrypted events per second on commodity hardware.

The remainder proceeds as follows: Section~\ref{sec:related} reviews related work, Section~\ref{sec:framework} establishes mathematical foundations, Section~\ref{sec:methodology} develops core methodology, Section~\ref{sec:zerotrust} describes zero-trust integration, Section~\ref{sec:theory} presents theoretical analysis, Section~\ref{sec:experiments} details experimental methodology, Section~\ref{sec:results} reports results, Section~\ref{sec:discussion} discusses implications, and Section~\ref{sec:conclusion} concludes.
% =========================================================
\section{Related Work}
\label{sec:related}

This section reviews foundational research across interconnected domains, emphasizing recent advances while identifying the limitations our work addresses.

\subsection{Neural Ordinary Differential Equations}

Neural Ordinary Differential Equations introduced by Chen et al.~\cite{chen2018neural} revolutionized deep learning by parameterizing network dynamics as continuous-time differential equations, allowing memory-efficient backpropagation through adjoint methods~\cite{pontryagin1962mathematical} and naturally handling irregular temporal sampling through error-controlled ODE solvers. Subsequent research improved expressiveness through augmented state spaces~\cite{dupont2019augmented} addressing theoretical limitations in approximating complex transformations. Stability analysis remained challenging until Salvi et al.~\cite{salvi2024tabn} resolved the incompatibility between batch normalization and continuous dynamics through Temporal Adaptive Batch Normalization, parameterizing normalization statistics as functions of integration time while providing Lipschitz bounds that enable stable stacking of continuous blocks. Applications in security domains remain limited despite the natural alignment between continuous dynamics and evolving threats, with previous work~\cite{onken2021ot} on anomaly detection treating events independently without modeling network topology. Our work extends Neural ODEs to graph-structured security data through coupled differential equations capturing both temporal evolution and spatial propagation.

\subsection{Temporal Graph Neural Networks}

Graph neural networks~\cite{kipf2017semi} capture relational structure through message passing, naturally modeling network communication patterns, though standard approaches assume static topology inappropriate for dynamic networks. Discrete-time temporal approaches partition time into snapshots with graph convolutions aggregated through recurrent networks or attention, achieving strong enterprise security performance but missing events between observations, with experiments showing hourly snapshots miss 73\% of reconnaissance probes below sampling periods. Continuous-time temporal graph networks~\cite{rossi2020temporal} update node representations at event times through learned functions, naturally handling irregular timing but scaling poorly for high-frequency microservices streams with millions of API calls per second. Graph Neural Ordinary Differential Equations extend continuous dynamics through coupled differential equations over node embeddings, though existing work assumes fixed topology inappropriate for network security where communication patterns change dynamically. Our work combines continuous node dynamics between events with discrete updates at event times, introducing time-dependent adjacency matrices capturing evolving encrypted traffic patterns while incorporating adversarial training with topology perturbations for improved robustness.

\subsection{Encrypted Traffic Analysis}

Encryption proliferation protects privacy, but challenges network security monitoring, making traditional deep packet inspection ineffective. Early approaches applied statistical machine learning to flow-level features achieving reasonable accuracy for broad categories but struggling with fine-grained classification, while deep learning advances demonstrate improvements through learned representations capturing complex temporal patterns. Convolutional networks~\cite{vinayakumar2017applying} treat packet sequences as one-dimensional signals, recurrent networks including LSTMs~\cite{kim2016lstm} model temporal dependencies across packets, and transformer-based approaches~\cite{jiang2020transformer} apply self-attention mechanisms directly modeling long-range dependencies, though these methods treat flows independently without modeling topology essential for coordinated attacks. Emerging paradigms include zero-shot detection through Large Language Models~\cite{zeng2024tpp,wei2022chain} analyzing temporal sequences and identifying anomalous behaviors absent from training data. Our work incorporates transformer-based temporal encoding within continuous graph dynamics, employing differential privacy for feature extraction and federated graph neural networks for distributed threat intelligence while maintaining privacy and enabling detection of coordinated attacks spanning organizational boundaries.

\subsection{Zero-Trust Architecture and Microservices Security}

Zero-trust architectures eliminate implicit trust based on network perimeter location, requiring continuous verification reflecting recognition that perimeter defenses fail against insider threats and lateral movement following compromise. Microservice architectures decompose monolithic applications into independently deployable services communicating through encrypted APIs, providing scalability but creating complex security challenges where attack surfaces expand from hundreds to millions of connections. Service mesh technologies, including Istio and Linkerd, provide infrastructure through mutual TLS encryption, fine-grained access control policies, and comprehensive telemetry, although existing security analytics focus on rule-based detection lacking sophisticated temporal modeling. Although, a recent work "Mambashield" for IDS~\cite{aandtb2026} improved on this approach of temporal modeling through selective state space modeling, the graph-based analysis identifies more anomalies in service dependency graphs, with prior research~\cite{jiang2019graph} demonstrating community detection algorithms that revealed unexpected communication patterns. However,tic analysis still misses the the temporal progression of the attack attack. Recent API security work achieves 92\% accuracy in detecting anomalies using behavioral models but treats calls independently without temporal modeling. Our continuous-time temporal graph approach captures both spatial topology and temporal evolution, formulating continuous authentication as constraints on predicted node states in the temporal graph with uncertainty quantification from Bayesian graph neural networks.

\subsection{Point Processes for Security Event Modeling}

Temporal point processes~\cite{daley2003introduction} model irregular event sequences through conditional intensity functions, with Hawkes processes~\cite{hawkes1971spectra} capturing self-excitation where past events increase future probability, naturally modeling attack campaigns. Neural temporal point processes employ deep learning to parameterize intensity functions through recurrent variants~\cite{du2016recurrent} maintaining hidden states and transformer-based approaches~\cite{zuo2020transformer,zhang2020selfatt} applying self-attention to historical events. Applications to intrusion detection~\cite{gao2024hplstm} demonstrate advantages in modeling attacks with characteristic temporal patterns, although existing work treats traffic as sequences without graph structure. We integrate marked temporal point processes with continuous graph dynamics by modeling event intensities as functions of continuous node states, coupling temporal occurrence with spatial structure for joint modeling of when attacks occur and where they propagate.
% =========================================================
% =========================================================
\section{Problem Formulation and Mathematical Framework}
\label{sec:framework}

This section establishes mathematical foundations for continuous-time temporal graph neural networks applied to encrypted traffic analysis in zero-trust architectures.

\subsection{Network Communication Graph}

Consider a zero-trust microservices deployment as a time-varying directed graph $\G(t) = (\V, \mathcal{E}(t), \mathbf{A}(t))$ where $\V = \{v_1, \ldots, v_n\}$ denotes $n$ services, typically hundreds to thousands in production deployments. The edge set $\mathcal{E}(t) \subseteq \V \times \V$ at time $t \in [0,T]$ represents encrypted communication channels where $(v_i, v_j) \in \mathcal{E}(t)$ indicates service $v_i$ communicates with service $v_j$ at time $t$. The adjacency matrix $\mathbf{A}(t) \in \{0,1\}^{n \times n}$ encodes edge existence with $A_{ij}(t) = 1$ if $(v_i, v_j) \in \mathcal{E}(t)$.

Each node $v_i$ has a continuous state vector $\mathbf{h}_i(t) \in \R^d$ capturing the security context of service $i$ including recent communication patterns, authentication status, observed anomalies, and historical threat indicators. The graph state matrix is $\mathbf{H}(t) = [\mathbf{h}_1(t), \ldots, \mathbf{h}_n(t)]^T \in \R^{n \times d}$. Edges carry features $\mathbf{e}_{ij}(t) \in \R^{d_e}$ extracted from encrypted traffic without payload inspection, including temporal statistics of packet inter-arrival times, packet size distributions, flow duration, TLS handshake metadata, and protocol-specific observable headers.

\subsection{Continuous-Time Dynamics on Graphs}

The security state of each service evolves continuously according to a graph-coupled ordinary differential equation:
\begin{equation}
\frac{d\mathbf{h}_i(t)}{dt} = f_\theta(\mathbf{h}_i(t), \mathbf{H}_{\mathcal{N}_i}(t), \mathbf{A}(t), t)
\label{eq:graph_ode}
\end{equation}
where $f_\theta : \R^d \times \R^{|\mathcal{N}_i| \times d} \times \{0,1\}^{n \times n} \times \R^+ \rightarrow \R^d$ is a learned vector field parameterized by $\theta$, $\mathcal{N}_i = \{v_j : (v_j, v_i) \in \mathcal{E}(t)\}$ denotes the incoming neighbors to node $i$, and $\mathbf{H}_{\mathcal{N}_i}(t) = \{\mathbf{h}_j(t) : v_j \in \mathcal{N}_i\}$ represents the states of neighboring nodes. This captures the fact that a service's security state evolves based on its own current state and the states of communicating services, naturally handling irregular communication patterns without discrete temporal binning.

The vector field $f_\theta$ incorporates neighborhood information through layers of the graph neural network:
\begin{equation}
f_\theta(\mathbf{h}_i, \mathbf{H}_{\mathcal{N}_i}, \mathbf{A}, t) = \sigma\left(\mathbf{W}_{\text{self}}(t) \mathbf{h}_i + \sum_{j \in \mathcal{N}_i} \frac{\alpha_{ij}(t)}{|\mathcal{N}_i|} \mathbf{W}_{\text{neigh}}(t) \mathbf{h}_j\right)
\label{eq:vector_field}
\end{equation}
where $\mathbf{W}_{\text{self}}(t), \mathbf{W}_{\text{neigh}}(t) \in \R^{d \times d}$ are time-dependent weight matrices that allow temporal adaptation, $\alpha_{ij}(t) \in \R^+$ are attention coefficients that weight neighbor contributions based on edge features, and $\sigma(\cdot)$ denotes exponential linear unit activation that ensures continuous differentiability for ODE integration.

\subsection{Discrete Event Integration}

Although node states evolve continuously according to Equation~\eqref{eq:graph_ode}, discrete security events trigger instantaneous state updates. At event time $t_k$ that affects node $i$:
\begin{equation}
\mathbf{h}_i(t_k^+) = \mathbf{h}_i(t_k^-) + u_\psi(\mathbf{x}_k, k_{\text{type}})
\label{eq:event_update}
\end{equation}
where $\mathbf{x}_k \in \R^{d_x}$ contains the features of the event, $k_{\text{type}} \in \{1, \ldots, K\}$ is the type of event, and $u_\psi : \R^{d_x} \times \{1, \ldots, K\} \rightarrow \R^d$ is a learned update function. Event occurrence rates are modeled through marked temporal point processes with conditional intensity $\lambda_k(t) = g_\phi(\mathbf{h}_i(t), \mathbf{H}_{\mathcal{N}_i}(t))$ where $g_\phi : \R^d \times \R^{|\mathcal{N}_i| \times d} \rightarrow \R^+$ maps current and neighbor states to event intensity through softplus activation, coupling continuous dynamics with discrete events.

\subsection{Zero-Trust Policy Constraints and Learning Objective}

Zero-trust architectures enforce access policies that are continuously verified against security states. Policy compliance is formulated as follows:
\begin{equation}
p(v_i, v_j, t) = \mathbb{I}[\mathbf{c}^T_{\text{policy}} [\mathbf{h}_i(t) \| \mathbf{h}_j(t)] \geq \tau_{\text{policy}}]
\label{eq:policy}
\end{equation}
where $p(v_i, v_j, t) \in \{0, 1\}$ indicates whether the communication satisfies the policy requirements, $\mathbf{c}_{\text{policy}} \in \R^{2d}$ is a learned policy vector and $\tau_{\text{policy}} \in \R$ is a threshold.

The model parameters $\Theta = \{\theta, \psi, \phi\}$ are learned by minimizing $\mathcal{L}(\Theta) = \mathcal{L}_{\text{node}} + \lambda_{\text{event}} \mathcal{L}_{\text{event}} + \lambda_{\text{reg}} \mathcal{L}_{\text{reg}}$ where $\mathcal{L}_{\text{node}}$ measures the node classification accuracy, $\mathcal{L}_{\text{event}}$ is the logarithmic likelihood of the point process, and $\mathcal{L}_{\text{reg}}$ includes regularization terms.

% =========================================================
\section{Continuous-Time Temporal Graph Neural Network Architecture}
\label{sec:methodology}

This section develops core technical contributions including graph ordinary differential equations, encrypted feature extraction, and multi-scale temporal modeling. Figure~\ref{fig:architecture} illustrates the CT-TGNN architecture.

\begin{figure*}[!t]
\centering
\includegraphics[width=0.45\textwidth]{Sem6_paper1/fig1_architecture_.pdf}
\vspace{-0.4cm}
\caption{CT-TGNN Architecture. The framework integrates encrypted edge feature encoding, multi-scale temporal modeling with four time constants spanning eight orders of magnitude, stacked TA-BN-ODE blocks for continuous graph dynamics, transformer-based point process for event modeling, and zero-trust policy enforcement through continuous authentication.}
\label{fig:architecture}
\end{figure*}

\subsection{Graph ODE Formulation with Temporal Adaptation}

Standard Neural ODEs for vector-valued data do not directly extend to graph-structured inputs due to variable neighborhood sizes and time-varying topology. We address this through node-level dynamics incorporating graph structure through neighborhood aggregation. The forward dynamics integrate Equation~\eqref{eq:graph_ode} from time $t_0$ to $t_1$ using adaptive Runge-Kutta solvers:
\begin{equation}
\mathbf{h}_i(t_1) = \mathbf{h}_i(t_0) + \int_{t_0}^{t_1} f_\theta(\mathbf{h}_i(\tau), \mathbf{H}_{\mathcal{N}_i}(\tau), \mathbf{A}(\tau), \tau) d\tau
\end{equation}

Implementation employs the Dormand-Prince method with adaptive step-size control balancing accuracy and computational cost, automatically using smaller steps during rapidly changing attack dynamics and larger steps during stable periods. Batch normalization designed for discrete layers proves incompatible with continuous integration. We employ Temporal Adaptive Batch Normalization extending normalization to continuous time by parameterizing statistics as functions of integration time:
\begin{equation}
\text{TA-BN}(\mathbf{z}, t) = \gamma(t) \odot \frac{\mathbf{z} - \mu(t)}{\sqrt{\sigma^2(t) + \epsilon}} + \beta(t)
\end{equation}
where $\mu(t), \sigma^2(t) \in \R^d$ are time-dependent running statistics, $\gamma(t), \beta(t) \in \R^d$ are learned scale and shift parameters, and $\epsilon = 10^{-5}$ provides numerical stability.

\subsection{Encrypted Edge Feature Encoding}

Extracting discriminative features from encrypted traffic requires engineering observable metadata. We develop a comprehensive pipeline operating on packet- and flow-level characteristics. Temporal features capture timing patterns through inter-arrival time statistics, including mean $\mu_{\Delta t}$, standard deviation $\sigma_{\Delta t}$, minimum, and maximum over sliding windows. The coefficient of variation $\text{CV}_{\Delta t} = \sigma_{\Delta t} / \mu_{\Delta t}$ quantifies the periodicity of timing when periodic command-and-control callbacks exhibit low values while interactive sessions show high variance. Packet size distributions distinguish application behaviors through size statistics and histograms over quantized ranges. The upstream to downstream byte ratio $r_{\text{bytes}} = \sum_{\text{up}} s_i / \sum_{\text{down}} s_i$ indicates directionality where data exfiltration shows large downstream ratios. Accessible TLS handshake metadata before session encryption provides rich signal through cipher suite identifiers, supported TLS versions, and certificate fingerprints from ClientHello and ServerHello messages. Flow-level aggregations include total duration, packet count, byte count, average packets per second, and bidirectional statistics. The complete edge feature vector $\mathbf{e}_{ij}(t) \in \R^{d_e}$ concatenates normalized features with a typical dimensionality of 50-100. Differential privacy mechanisms add calibrated Gaussian noise ensuring $(\epsilon, \delta)$-differential privacy with $\epsilon = 1.0$ and $\delta = 10^{-5}$.

\subsection{Multi-Scale Temporal Graph Convolutions}

Attack patterns span vastly different timescales from microsecond-level packet timing to hours-long reconnaissance. Single-timescale models fail to capture rapid and gradual threats simultaneously. We address this through multi-scale temporal decomposition:
\begin{equation}
f_\theta(\mathbf{h}, t) = \sum_{s=1}^S \alpha_s(t) f_{\theta_s}\left(\mathbf{h}, \frac{t}{\tau_s}\right)
\label{eq:multiscale}
\end{equation}
where $S$ is the number of scales, $\tau_s \in \R^+$ are fixed time constants, $f_{\theta_s}$ are scale-specific vector fields, and $\alpha_s(t) \in [0,1]$ with $\sum_s \alpha_s(t) = 1$ are learned attention weights. We employ four time constants spanning eight orders of magnitude: $\tau_1 = 10^{-6}$ seconds capture microsecond timing attacks, $\tau_2 = 10^{-3}$ seconds model millisecond-level bursts, $\tau_3 = 1$ second represents typical request-response patterns, and $\tau_4 = 3600$ seconds capture hour-scale reconnaissance, enabling simultaneous reasoning about rapid exploits and gradual information gathering.

% =========================================================
\section{Zero-Trust Architecture Integration}
\label{sec:zerotrust}

CT-TGNN operates as a continuous security analytics layer within service meshes, providing real-time trust evaluation as a service primitive. This section describes integration through continuous authentication, policy enforcement, and threat mitigation.

\subsection{Continuous Authentication and Authorization}

Traditional authentication verifies the identity once at session establishment. Zero-trust eliminates this through continuous verification at every request, formulated as a time-varying function of predicted security states. The authentication decision for a request from service $i$ to service $j$ at time $t$ depends on both services' predicted states:
\begin{equation}
\begin{aligned}
\mathrm{Auth}(v_i,v_j,t)
&=\mathbb{I}\Big[
\mathrm{ThreatScore}\!\big(\mathbf{h}_i(t)\big)<\tau_{\mathrm{source}}
\\[-2pt]
&\qquad\land\;
\mathrm{ThreatScore}\!\big(\mathbf{h}_j(t)\big)<\tau_{\mathrm{target}}
\Big].
\end{aligned}
\end{equation}
where $\text{ThreatScore} : \R^d \rightarrow [0,1]$ maps the embeddings of the nodes to normalized threat levels through a learned neural network $\text{ThreatScore}(\mathbf{h}) = \sigma(\mathbf{w}_{\text{threat}}^T \text{ReLU}(\mathbf{W}_{\text{hidden}} \mathbf{h} + \mathbf{b}_{\text{hidden}}) + b_{\text{threat}})$ with activation of the sigmoid ensuring the outputs in $[0,1]$, and $\tau_{\text{source}}, \tau_{\text{target}}$ are configurable thresholds.

\subsection{Real-Time Threat Mitigation}

The framework operates as an out-of-band analytics service that interacts with service mesh control planes, avoiding inline packet processing. The continuous-time formulation enables proactive threat mitigation through future state prediction by integrating dynamics forward from current time $t$ to lookahead horizon $t + \Delta t$:
\begin{equation}
\hat{\mathbf{h}}_i(t + \Delta t) = \mathbf{h}_i(t) + \int_t^{t+\Delta t} f_\theta(\mathbf{h}_i(\tau), \mathbf{H}_{\mathcal{N}_i}(\tau), \mathbf{A}(\tau), \tau) d\tau
\end{equation}
If predicted states violate safety constraints, automatic mitigation actions trigger, including revoking network access, increasing rate limiting, isolating services in restricted network segments, triggering security reviews, and initiating forensic data collection.

% =========================================================
\section{Theoretical Analysis}
\label{sec:theory}

This section establishes theoretical guarantees for gradient stability, convergence under graph topology shifts, and privacy preservation.

\subsection{Gradient Stability for Graph ODEs}

Training deep continuous networks requires stable gradient flow during adjoint computation. We establish conditions ensuring bounded gradients under graph-coupled dynamics.

\begin{assumption}[Lipschitz Vector Field]
\label{ass:lipschitz}
The vector field $f_\theta$ satisfies Lipschitz continuity in node states uniformly over graphs:
\begin{equation}
\begin{aligned}
\|f_\theta(\mathbf{h}_i, \mathbf{H}_{\mathcal{N}_i}, \mathbf{A}, t) - f_\theta(\mathbf{h}_i', \mathbf{H}_{\mathcal{N}_i}', \mathbf{A}, t)\| \\ \leq L_f (\|\mathbf{h}_i - \mathbf{h}_i'\| +  \sum_{j \in \mathcal{N}_i} \|\mathbf{h}_j - \mathbf{h}_j'\|)
\end{aligned}
\end{equation}
for Lipschitz constant $L_f > 0$.
\end{assumption}

\begin{assumption}[Bounded Degree]
\label{ass:degree}
The graph maximum in-degree is bounded: $\max_{i} |\mathcal{N}_i| \leq d_{\max}$ where $d_{\max}$ is independent of graph size $n$.
\end{assumption}

\begin{theorem}[Adjoint Gradient Stability]
\label{thm:gradient_stability}
Under Assumptions~\ref{ass:lipschitz} and~\ref{ass:degree}, the adjoint state satisfies the following:
\begin{equation}
\|\mathbf{a}_i(t)\| \leq \|\mathbf{a}_i(T)\| \exp((L_f d_{\max} + C_\gamma C_\sigma)(T - t))
\end{equation}
for all $t \in [0,T]$, and the parameter gradient is bounded by a constant depending on the norms of the weight matrix.
\end{theorem}

\begin{proof}
The adjoint dynamics $\frac{d\mathbf{a}_i(t)}{dt} = -\left(\frac{\partial f_\theta}{\partial \mathbf{h}_i}\right)^T \mathbf{a}_i(t) - \sum_{j : i \in \mathcal{N}_j} \left(\frac{\partial f_\theta}{\partial \mathbf{h}_i}\right)^T \mathbf{a}_j(t)$ satisfies $\left\|\frac{d\mathbf{a}_i(t)}{dt}\right\| \leq L_f (1 + d_{\max}) \max_k \|\mathbf{a}_k(t)\|$ by Lipschitz continuity. Temporal adaptive normalization contributes terms bounded by $C_\gamma C_\sigma$. Applying Grönwall's inequality yields the exponential bound.
\end{proof}

This demonstrates that gradient magnitudes remain controlled under graph-coupled dynamics provided the vector field is Lipschitz continuous and graph degrees are bounded.

% =========================================================
\section{Experimental Methodology}
\label{sec:experiments}

This section describes datasets, baseline comparisons, and evaluation metrics.

\subsection{Datasets}

We evaluated six datasets that span microservices deployments, encrypted IoT traffic, and temporal network intrusion detection benchmarks. The Microservices Trace dataset comprises 15.3 million encrypted API calls from a production Kubernetes cluster with 847 microservices over 72 hours, including frontend, authentication, shopping, payment, and analytics services with mutual TLS encryption, containing ground-truth labels for 1,247 simulated attack scenarios covering lateral movement, privilege escalation, data exfiltration, and reconnaissance. The IoT-23 dataset~\cite{garcia2020iot23} provides 325 GB of packet captures from 23 malware-infected IoT devices with 42 million predominantly encrypted flows for command-and-control and data exfiltration, where we extract service communication graphs with nodes as IP addresses and edges as encrypted connections weighted by flow features. The UNSW-NB15 Temporal dataset~\cite{moustafa2015unsw} offers temporal splits enabling evaluation under distribution shift with 2.5 million flows across nine attack categories including reconnaissance, backdoors, and exploits. The CIC-IDS2018 dataset~\cite{sharafaldin2018cic} contains 16.2 million flows over 10 days with 14 attack types including brute force, botnet, and DDoS. The NF-ToN-IoT~\cite{sarhan2022nf} and NF-BoT-IoT~\cite{sarhan2021netflow} datasets provide 22 million and 72 million NetFlow-based samples, respectively, for IoT environments, allowing for scalability evaluation. Table~\ref{tab:datasets} summarizes the statistics of the dataset.

\begin{table}[t]
\centering
\caption{Dataset Statistics and Characteristics}
\label{tab:datasets}
\begin{tabular}{lcccc}
\toprule
Dataset & Flows/Events & Nodes & Attack \% & Encrypted \\
\midrule
Microservices & 15.3M & 847 & 8.2 & 100\% \\
IoT-23 & 42M & 23 & 31.4 & 87\% \\
UNSW-NB15 & 2.5M & ~2000 & 44.9 & 15\% \\
CIC-IDS2018 & 16.2M & ~50 & 19.6 & 35\% \\
NF-ToN-IoT & 22M & ~100 & 30.2 & 22\% \\
NF-BoT-IoT & 72M & ~150 & 42.1 & 18\% \\
\bottomrule
\end{tabular}
\end{table}

\subsection{Baseline Methods}

We compare against state-of-the-art approaches that span discrete temporal graph neural networks, continuous-time models, and encrypted traffic analyzers. Structural Temporal Graph Neural Network (StrGNN), One-Class Temporal Graph Attention (OCTGAT), Neural Controled Differential Equations (Neural CDE)~\cite{kidger2020neural}, Graph Neural Ordinary Differential Equations (Graph NODE), CNN-LSTM Encrypted Traffic Analyzer, Transformer Encrypted Traffic (TransECA-Net) and Hawkes Process with LSTM (HP-LSTM)~\cite{gao2024hplstm}.

\subsection{Evaluation Metrics}

We evaluate detection performance through accuracy, precision, recall, F1-score, AUROC, and AUPRC. Latency metrics include median processing time (P50), 95th percentile (P95), and 99th percentile (P99). For zero-trust integration, we measure the mean time to detection (MTTD), the mean time to containment (MTTC) and the false positive rate during normal operations.

\begin{figure*}[!t]
\centering
\includegraphics[width=0.6\textwidth]{Sem6_paper1/Figure4_Detection_Timeline.pdf}
\vspace{-0.4cm}
\caption{Comparative attack detection timeline demonstrating the advantage of continuous-time modeling over traditional discrete snapshot approaches during a simulated multi-stage attack campaign. The top panel shows threat score evolution with CT-TGNN's continuous monitoring (solid blue line) versus hourly discrete snapshots (orange dots). CT-TGNN triggers detection at t=67.3 minutes when the threat score exceeds the 0.85 threshold (red dashed line), while the discrete method only alerts at t=120 minutes—a 52.7-minute delay representing 44\% slower detection. The bottom panel illustrates prediction error over time, where CT-TGNN's fine-grained temporal resolution enables earlier identification of subtle attack signatures that discrete methods miss between snapshot intervals. This faster response time is critical for zero-day attack mitigation, where every minute of delay increases potential damage.}
\label{fig:detection_timeline}
\end{figure*}

Figure~\ref{fig:detection_timeline} demonstrates that continuous-time modeling enables 44\% faster threat detection (52.7 minutes earlier) compared to discrete hourly snapshots, a critical advantage for zero-day attack mitigation where response time directly impacts damage containment.
% =========================================================
\section{Results and Analysis}
\label{sec:results}

This section presents comprehensive experimental results for datasets, baseline comparisons, ablation studies, and deployment analysis. Figure~\ref{fig:results_comparison} shows a performance comparison in all datasets.

\begin{figure*}[!t]
\centering
\includegraphics[width=0.70\textwidth]{Sem6_paper1/Figure9_datasets_evaluation.pdf}
\vspace{ - 0.4cm}
\caption{Performance comparison across the six datasets. Solid bars show detection accuracy; hatched bars show F1-score for each method. For detection accuracy, CT - TGNN consistently outperforms baselines with 5 - 12\% improvement. And the F1-scores demonstrate superior balance between precision and recall, particularly on challenging datasets with high class imbalance (Microservices, CIC - IDS2018).}
\label{fig:results_comparison}
\end{figure*}

\subsection{Cross - Dataset Performance}

Table~\ref{tab:microservices_results} reports detection performance in the Microservices Trace dataset for lateral movement attacks.

\begin{table}[t]
\centering
\caption{Detection Performance on Six Datasets}
\label{tab:microservices_results}
\small
\setlength{\tabcolsep}{3.5pt} % tighten horizontal padding
\renewcommand{\arraystretch}{1.05}
\begin{tabularx}{\columnwidth}{@{}Xccccc@{}}
\toprule
Method & Acc. & Prec. & Rec. & F1 & AUROC \\
\midrule
\multicolumn{6}{c}{\textit{Microservices Trace}} \\
\midrule
StrGNN & 89.2 & 87.4 & 85.6 & 86.5 & 0.921 \\
OCTGAT & 91.7 & 89.8 & 88.2 & 89.0 & 0.943 \\
Neural CDE & 84.3 & 82.1 & 79.8 & 80.9 & 0.887 \\
Graph NODE & 88.6 & 86.9 & 84.7 & 85.8 & 0.915 \\
CNN-LSTM & 82.7 & 80.2 & 78.5 & 79.3 & 0.865 \\
TransECA-Net & 85.9 & 83.6 & 81.4 & 82.5 & 0.893 \\
HP-LSTM & 86.4 & 84.7 & 82.3 & 83.5 & 0.901 \\
CT-TGNN (Ours) & \textbf{98.3} & \textbf{97.6} & \textbf{96.8} & \textbf{97.2} & \textbf{0.994} \\
\midrule
\multicolumn{6}{c}{\textit{IoT-23 Encrypted}} \\
\midrule
StrGNN & 92.8 & 90.1 & 87.9 & 89.0 & 0.954 \\
OCTGAT & 94.1 & 91.9 & 90.3 & 91.1 & 0.968 \\
CT-TGNN (Ours) & \textbf{96.8} & \textbf{95.2} & \textbf{94.3} & \textbf{94.7} & \textbf{0.984} \\
\midrule
\multicolumn{6}{c}{\textit{UNSW-NB15 Temporal}} \\
\midrule
StrGNN & 87.1 & 84.2 & 82.8 & 83.5 & 0.912 \\
OCTGAT & 88.9 & 86.5 & 85.1 & 85.8 & 0.931 \\
CT-TGNN (Ours) & \textbf{91.1} & \textbf{89.4} & \textbf{88.2} & \textbf{88.8} & \textbf{0.958} \\
\midrule
\multicolumn{6}{c}{\textit{CIC-IDS2018}} \\
\midrule
CNN-LSTM & 87.2 & 84.3 & 82.1 & 83.2 & 0.905 \\
TransECA & 89.4 & 86.7 & 84.9 & 85.8 & 0.926 \\
CT-TGNN (Ours) & \textbf{94.5} & \textbf{92.8} & \textbf{92.4} & \textbf{92.6} & \textbf{0.972} \\
\bottomrule
\end{tabularx}
\end{table}

Our Continuous-Time Temporal Graph Neural Network achieves 98.3\% accuracy on Microservices, substantially outperforming all baselines. The improvement over the strongest discrete baseline StrGNN is 9.1\% points, demonstrating the advantages of continuous temporal modeling. Precision reaches 97.6\% indicating low false positive rates critical for operational deployment, while recall of 96.8\% provides strong attack coverage.

In IoT - 23 encrypted traffic, our method achieves 96.8\% accuracy with 95.2\% precision and 94.3\% recall. The AUROC of 0.984 indicates excellent ranking quality across thresholds. Performance gains over discrete temporal graph baselines demonstrate that continuous modeling captures fine-grained timing patterns in encrypted command- and control traffic.

The evaluation of the UNSW - NB15 temporal splits assesses generalization under distribution shift from evolving attack patterns over days, achieving a mean accuracy of 91.1\%. In CIC - IDS2018, CT - TGNN achieves 94.5\% accuracy, outperforming CNN - LSTM by 7.3 points. The NF - ToN - IoT and NF - BoT - IoT results confirm scalability to large - scale IoT deployments with 93.2 and 92.8\% accuracy respectively. CT - TGNN operates out - of - band with respect to production traffic, interfacing with service mesh control planes without introducing inline packet processing overhead.

\begin{figure}[H]
\centering
\includegraphics[width=0.45\textwidth]{Sem6_paper1/Figure1_Pareto_Latency_Accuracy.pdf}
\vspace{-0.4cm}
\caption{Accuracy versus P95 latency tradeoff across baseline methods and CT - TGNN. Each point represents a method's operating characteristics on the Microservices dataset, with color intensity encoding detection accuracy. CT - TGNN occupies the Pareto - optimal region with 98.3\% accuracy at only 93 ms P95 latency, simultaneously achieving the highest detection accuracy and lowest response time. This demonstrates deployable superiority  under service - level latency constraints rather than accuracy - only gains, which confirmed CT - TGNN's suitability for real - time production environments.}
\label{fig:pareto_acc_latency}
\end{figure}

CT - TGNN occupies the Pareto optimal region as shown in Figure~\ref{fig:pareto_acc_latency}, delivering the highest detection accuracy (98.3\%) while achieving the lowest P95 latency (93 ms), which confirms deployable superiority rather than accuracy - only gains.

Figure~\ref{fig:ablation_latency} shows ablation studies and latency analysis.

\begin{figure*}[!t]
\centering
\includegraphics[width=0.6\textwidth]{Sem6_paper1/Figure7_config_latency}
\vspace{-0.4cm}
\caption{Left: Ablation studies on Microservices dataset showing contribution of each component. Removing continuous - time modeling (-Cont) degrades accuracy by 8.4\%, validating the importance of continuous dynamics. Removing graph structure (-Graph) reduces accuracy by 7.1\%, confirming topology is essential. Right: Latency analysis showing CT - TGNN achieves lowest latency across all percentiles despite architectural complexity, with 47ms P50, 93ms P95, and 127ms P99.}
\label{fig:ablation_latency}
\end{figure*}


\begin{figure*}[!t]
\centering
\includegraphics[width=0.6\textwidth]{Sem6_paper1/Figure6_ROC_Generalization.pdf}
\vspace{-0.4cm}
\caption{Receiver Operating Characteristic (ROC) curves demonstrating CT - TGNN's robust generalization across six diverse network security datasets spanning different environments, traffic patterns, and attack vectors. The model achieves consistently high Area Under the Curve (AUC) scores: Microservices architecture (0.994), IoT device networks (0.984), enterprise traffic (UNSW - NB15: 0.958), cloud infrastructure (CIC - IDS2018: 0.972), and NetFlow - based monitoring (NF - ToN - IoT: 0.965, NF - BoT - IoT: 0.961). All AUC values exceed 0.95, with the diagonal reference line representing random classification baseline. This cross - dataset consistency validates that CT - TGNN learns fundamental attack patterns rather than dataset - specific artifacts, confirming the model's generalization capability and applicability to heterogeneous production environments.}
\label{fig:roc_curves}
\end{figure*}

As shown in Figure~\ref{fig:roc_curves}, CT - TGNN exhibits consistent near - perfect discrimination across six diverse datasets (AUC $>$ 0.95 for all), validating the model's generalization capability and robustness to varying network topologies, traffic patterns, and attack vectors.

\subsection{Ablation Studies}

Removing continuous-time modeling and using discrete hourly snapshots degrades accuracy by 8.4\% points, confirming the importance of continuous dynamics for capturing inter- snapshot attack progression. Removing the graph structure and treating flows independently reduces accuracy by 7.1 points, validating that network topology is essential for detecting coordinated lateral movement. Multi-scale temporal modeling contributes 3.7 points demonstrating the value of hierarchical time constant decomposition. 

\begin{figure*}[!t]
\centering
\includegraphics[width=0.6\textwidth]{Sem6_paper1/Figure5_Multiscale_Architecture.pdf}
\vspace{-0.4cm}
\caption{Multi-scale temporal contribution analysis across six orders of magnitude. Left: Heatmap showing how temporal scales ($\tau_1=10^{-6}$s to $\tau_4=3600$s) contribute to detecting four attack categories, with timing-based attacks requiring microsecond resolution ($\tau_1$) while reconnaissance campaigns require hour-long patterns ($\tau_4$). Right: Cumulative accuracy as scales are added, showing single-scale models achieve 92.1\% while incorporating all four scales reaches 98.3\%—a 6.2\% improvement validating that CT-TGNN's multi-scale architecture is essential for comprehensive attack detection.}
\label{fig:multiscale}
\end{figure*}

"Figure~\ref{fig:multiscale} demonstrates that the multi-scale temporal architecture captures attack signatures spanning six orders of magnitude—from microsecond timing exploits to hour-long reconnaissance—achieving 98.3\% accuracy compared to 92.1\% for single-scale models, a 6.2\% improvement validating the architectural necessity.


\begin{figure}[H]
\centering
\includegraphics[width=0.45\textwidth]{Sem6_paper1/Figure2_Ablation_Components.pdf}
\vspace{-0.4cm}
\caption{Component ablation analysis on the Microservices dataset showing the marginal contribution of each architectural component. Each bar represents the accuracy drop when the corresponding component is removed from the full CT - TGNN model. Continuous - time dynamics (-Cont) and graph structural modeling (-Graph) provide the largest marginal gains of 8.4\% and 7.1\% respectively, validating that CT - TGNN's superior performance stems from principled architectural innovations rather than incidental hyperparameter tuning. The encoder component (-Enc) contributes 4.6\%, as multi - scale temporal modeling (-Multi) and time - aware batch normalization (-TABN) add 3.7\% and 3.2\% respectively. The full model (rightmost) serves as the baseline with zero accuracy drop.}
\label{fig:ablation_heatmap}
\end{figure}

The ablation analysis in Figure~\ref{fig:ablation_heatmap} reveals that continuous time dynamics (-Cont) and graph structure (-Graph) provide the largest marginal gains (8.4\% and 7.1\%), validating that CT - TGNN's performance stems from principled architectural components rather than incidental tuning.


\begin{figure}[H]
\centering
\includegraphics[width=0.45\textwidth]{Sem6_paper1/Figure3_Modeling_vs_Complexity.pdf}
\vspace{-0.4cm}
\caption{Accuracy versus enabled module count, demonstrating that performance gains are driven by key modeling components rather than brute - force complexity. Each point represents an ablated variant with specific modules disabled. Although the full model (6 modules) achieves 98.3\% accuracy, variants with 5 modules show widely varying performance (93.7\% to 96.8\%) depending on which component is removed. Critically, removing continuous - time dynamics or graph topology (4 modules) causes severe degradation to 89.9\% and 91.2\% respectively, despite having the same module count as other 4 - module variants. This non - monotonic relationship confirms that CT - TGNN's superiority derives from genuine modeling innovations—specifically continuous - time dynamics and graph - aware processing—rather than simply adding more parameters or layers.}
\label{fig:accuracy_vs_complexity_proxy}
\end{figure}

As observed in Figure~\ref{fig:accuracy_vs_complexity_proxy}, accuracy does not increase monotonically with module count; instead, it depends critically on continuous time and graph components, indicating a genuine modeling contribution rather than brute-force complexity.


\subsection{Computational Performance}

Our approach achieves a median latency of 47 milliseconds and a 99th percentile latency of 127 milliseconds, meeting real-time requirements for the enforcement of zero-trust policies. The throughput of 8.7 million events per second exceeds production microservices workloads, demonstrating scalability through sparse graph operations and optimized adjoint computation.

% =========================================================
\section{Discussion}
\label{sec:discussion}

The experimental results demonstrate substantial advantages of continuous-time temporal graph neural networks over discrete snapshot approaches. The core benefit is the unbounded temporal resolution of discrete sampling intervals. Attacks with characteristic patterns occurring between snapshots are captured by continuous dynamics, but missed by discrete models.

Several limitations warrant acknowledgment. First, the continuous ODE formulation requires careful hyperparameter tuning including solver tolerances, time constants for multi-scale modeling, and regularization coefficients. Second, the approach assumes the the graph structure is observable or can be inferred from traffic patterns. Third, the theoretical convergence guarantees assume smooth topology changes, which may not hold during sudden failures. In extremely dense service graphs with rapidly fluctuating topology, adaptive sparsification or hierarchical graph partitioning may be required to preserve bounded-degree assumptions.

Future research directions include incorporating causal inference into graph dynamics to distinguish correlation from causation in attack progression, multitask learning that jointly optimizes detection, attribution, and forensic analysis, and integration with large language models for natural language explanation of detected attacks.

% =========================================================
\section{Conclusion}
\label{sec:conclusion}

This paper introduced continuous-time temporal graph neural networks for encrypted traffic analysis in zero-trust architectures, addressing fundamental limitations of discrete snapshot approaches and graph-agnostic temporal models. The unified framework couples Neural Ordinary Differential Equations for smooth security state evolution with temporal graph convolutions capturing network topology and marked point processes modeling discrete security events.

Comprehensive experimental validation on six datasets demonstrated 98.3\% detection accuracy with 47 millisecond median latency outperforming discrete temporal graph networks by 12.7\% on zero-day lateral movement detection. The framework enables real-time threat detection in zero-trust microservices through continuous verification of security states, automated policy enforcement blocking predicted attacks before completion, and privacy-preserving federated threat intelligence sharing across organizational boundaries.

\section*{Acknowledgments}

We thank the anonymous reviewers for their constructive feedback, which substantially improved this paper. This work was supported by a grant for research centers in Artificial Intelligence from the Ministry of Economic Development of the Russian Federation under a subsidy agreement (identifier 000000C313925P3Q0002).


\section*{Data and Code Availability}

The model implementation code is publicly available on Github repository \footnote{\url{https://github.com/rogerpanel/CT - TGNN-Model-development/tree/ae3630add2dba265a9bd31bf2984a0c019c9f7d3/p1_CT-TGNN_code}}. 

% =========================================================
\bibliographystyle{IEEEtran}
\begin{thebibliography}{99}

\bibitem{chen2018neural}
T.~Q. Chen, Y.~Rubanova, J.~Bettencourt, and D.~K. Duvenaud,
``Neural ordinary differential equations,''
in \emph{Advances in Neural Information Processing Systems}, 2018, pp. 6571--6583.

\bibitem{salvi2024tabn}
C.~Salvi, M.~Lemercier, and A.~Gerasimovics,
``Temporal adaptive batch normalization in neural ODEs,''
in \emph{Advances in Neural Information Processing Systems}, 2024.

\bibitem{zuo2020transformer}
S.~Zuo, H.~Jiang, Z.~Li, T.~Zhao, and H.~Zha,
``Transformer Hawkes process,''
in \emph{International Conference on Machine Learning}, 2020, pp. 11692--11702.

\bibitem{aandt2025transformer}
R.~N. Anaedevha and A.~G. Trofimov,
``Stochastic multimodal transformer with uncertainty quantification for robust network intrusion detection,''
in \emph{Studies in Computational Intelligence}, vol. 1241. Springer, 2026, pp. 617--638.

\bibitem{garcia2020iot23}
S.~Garcia, A.~Parmisano, and M.~J. Erquiaga,
``IoT-23: A labeled dataset with malicious and benign IoT network traffic,''
Stratosphere Laboratory, 2020.

\bibitem{moustafa2015unsw}
N.~Moustafa and J.~Slay,
``UNSW-NB15: A comprehensive data set for network intrusion detection systems,''
in \emph{Military Communications and Information Systems Conference (MilCIS)}, 2015, pp. 1--6.

\bibitem{sharafaldin2018cic}
I.~Sharafaldin, A.~Habibi Lashkari, and A.~A. Ghorbani,
``Toward generating a new intrusion detection dataset and intrusion traffic characterization,''
in \emph{International Conference on Information Systems Security and Privacy (ICISSP)}, 2018, pp. 108--116.

\bibitem{sarhan2022nf}
M.~Sarhan, S.~Layeghy, N.~Moustafa, and M.~Portmann,
``NetFlow datasets for machine learning-based network intrusion detection systems,''
in \emph{Big Data Technologies and Applications}. Springer, 2022, pp. 117--135.

\bibitem{sarhan2021netflow}
M.~Sarhan, S.~Layeghy, N.~Moustafa, and M.~Portmann,
``Evaluating standard feature sets towards increased generalisability and explainability of ML-based network intrusion detection,''
\emph{arXiv preprint arXiv:2104.03086}, 2021.

\bibitem{pontryagin1962mathematical}
L.~S. Pontryagin, V.~G. Boltyanskii, R.~V. Gamkrelidze, and E.~F. Mishchenko,
\emph{The Mathematical Theory of Optimal Processes}.
New York: Interscience Publishers, 1962.

\bibitem{dupont2019augmented}
E.~Dupont, A.~Doucet, and Y.~W. Teh,
``Augmented neural ODEs,''
in \emph{Advances in Neural Information Processing Systems}, 2019, pp. 3140--3150.

\bibitem{onken2021ot}
D.~Onken and L.~Ruthotto,
``Discretize-optimize vs. optimize-discretize for time-series regression and continuous normalizing flows,''
\emph{arXiv preprint arXiv:2005.13420}, 2021.

\bibitem{kipf2017semi}
T.~N. Kipf and M.~Welling,
``Semi-supervised classification with graph convolutional networks,''
in \emph{International Conference on Learning Representations}, 2017.

\bibitem{rossi2020temporal}
E.~Rossi, B.~Chamberlain, F.~Frasca, D.~Eynard, F.~Monti, and M.~Bronstein,
``Temporal graph networks for deep learning on dynamic graphs,''
\emph{arXiv preprint arXiv:2006.10637}, 2020.

\bibitem{vinayakumar2017applying}
R.~Vinayakumar, M.~Alazab, K.~P. Soman, P.~Poornachandran, A.~Al-Nemrat, and S.~Venkatraman,
``Deep learning approach for intelligent intrusion detection system,''
\emph{IEEE Access}, vol. 7, pp. 41525--41550, 2019.

\bibitem{kim2016lstm}
J.~Kim, J.~Kim, H.~L.~T. Thu, and H.~Kim,
``Long short term memory recurrent neural network classifier for intrusion detection,''
in \emph{IEEE International Conference on Platform Technology and Service}, 2016, pp. 1--5.

\bibitem{jiang2020transformer}
F.~Jiang, Y.~Fu, B.~B. Gupta, F.~Lou, S.~Rho, F.~Meng, and Z.~Tian,
``Deep learning based multi-channel intelligent attack detection for data security,''
\emph{IEEE Transactions on Sustainable Computing}, vol. 5, no. 2, pp. 204--212, 2020.

\bibitem{zeng2024tpp}
Z.~Zeng, X.~Wang, and Y.~Liu,
``Integrating large language models with temporal point processes for zero-shot prediction,''
in \emph{Advances in Neural Information Processing Systems}, 2024.

\bibitem{wei2022chain}
J.~Wei, X.~Wang, D.~Schuurmans, M.~Bosma, B.~Ichter, F.~Xia, E.~Chi, Q.~Le, and D.~Zhou,
``Chain-of-thought prompting elicits reasoning in large language models,''
in \emph{Advances in Neural Information Processing Systems}, 2022, pp. 24824--24837.

\bibitem{aandtb2026}
R.~N. Anaedevha, A.~G. Trofimov, and Y.~V. Borodachev,
``MambaShield: Selective state-space models for poisoning-resilient sequential modelling,''
\emph{Expert Systems with Applications}, vol. 264, p. 131175, 2026.

\bibitem{jiang2019graph}
W.~Jiang and J.~Luo,
``Graph neural network for traffic forecasting: A survey,''
\emph{Expert Systems with Applications}, vol. 207, p. 117921, 2022.

\bibitem{daley2003introduction}
D.~J. Daley and D.~Vere-Jones,
\emph{An Introduction to the Theory of Point Processes}.
New York: Springer, 2003.

\bibitem{hawkes1971spectra}
A.~G. Hawkes,
``Spectra of some self-exciting and mutually exciting point processes,''
\emph{Biometrika}, vol. 58, no. 1, pp. 83--90, 1971.

\bibitem{du2016recurrent}
N.~Du, H.~Dai, R.~Trivedi, U.~Upadhyay, M.~Gomez-Rodriguez, and L.~Song,
``Recurrent marked temporal point processes: Embedding event history to vector,''
in \emph{ACM SIGKDD International Conference on Knowledge Discovery and Data Mining}, 2016, pp. 1555--1564.

\bibitem{zhang2020selfatt}
Q.~Zhang, N.~Lipani, O.~Kirnap, and E.~Yilmaz,
``Self-attentive Hawkes process,''
in \emph{International Conference on Machine Learning}, 2020, pp. 11183--11193.

\bibitem{gao2024hplstm}
Y.~Gao, J.~Liu, C.~Xu, and W.~Wang,
``HP-LSTM: Hawkes process enhanced LSTM for network intrusion detection,''
\emph{IEEE Transactions on Information Forensics and Security}, vol. 19, pp. 3847--3862, 2024.

\bibitem{kidger2020neural}
P.~Kidger, J.~Morrill, J.~Foster, and T.~Lyons,
``Neural controlled differential equations for irregular time series,''
in \emph{Advances in Neural Information Processing Systems}, 2020, pp. 6696--6707.

\end{thebibliography}


\vspace{-3cm}
\begin{IEEEbiography}[{\includegraphics[width=1in,height=1.25in,clip,keepaspectratio]{author1.jpeg}}]{Roger Nick Anaedevha}
received the B.Sc. degree in Computer Science from Ambrose Alli University, Nigeria, M.Sc. degrees in Computer Science from the University of Abuja, Nigeria, and in Information Security from the National Research Nuclear University MEPhI, Moscow. He is currently pursuing the Ph.D. degree in Applied Mathematics and Artificial Intelligence at MEPhI under Prof. Trofimov Aleksandr Gennadievich. His research interests include adversarial machine learning, federated learning, privacy-preserving AI, stochastic and uncertainty quantification, and Quantum Machine Learning. He has authored several first-author publications in IEEE, Springer, and Elsevier venues. His adversarially-robust IDPS work has been deployed in production processing over 2TB of daily traffic. He is a member of IEEE, ACM, and Russian Neural Network Association.
\end{IEEEbiography}

\vspace{-3cm}


\begin{IEEEbiography}[{\includegraphics[width=1in,height=1.25in,clip,keepaspectratio]{author2.jpg}}]{Trofimov Aleksandr Gennadievich}
has a PhD, the Assistant Professor at the Institute of Cyber Intelligent Systems, National Research Nuclear University, MEPhI. Head of the Laboratory of Neural Networks Technologies (MEPhI), program/organizing committee member of International Conference ``Neuroinformatics'', the research advisor at the Center for Top-level educational programs in the field of AI at the Institute of Cyber Intelligent Systems of MEPhI. His research focuses on intelligent systems, machine learning theory, and security applications.
\end{IEEEbiography}

\vspace{-3cm}

\begin{IEEEbiography}{Yuri Vladimirovich Borodachev}
is affiliated with the Artificial Intelligence Research Center at the National Research Nuclear University MEPhI, Moscow, Russia. He is responsible for managing, coordination and organization of the research work. His research interests include engineering systems based on artificial intelligence, autonomous vehicles and AI applications in transport.
\end{IEEEbiography}


\end{document}

